\documentclass[a4paper,12pt]{article}
\usepackage[utf8]{inputenc}
\usepackage[T2A]{fontenc}
\usepackage[russian]{babel}
\usepackage{amsmath}
\usepackage{amssymb}
\usepackage{array}
\usepackage{booktabs}
\usepackage{geometry}
\usepackage{xcolor}
\usepackage{tabularx}
\usepackage{multirow}

\geometry{top=2cm, bottom=2cm, left=2.5cm, right=2cm}

\title{\textbf{Домашняя работа 5 по дискретной математике}\\[0.5em]
\large Изоморфизм графов. Вычисление кода графа (Раздел~9.2)}
\author{Кливцов Александр Александрович \\ ISU: 504590, Вариант 29}
\date{}

\begin{document}

\maketitle
\thispagestyle{empty}

% ──────────────────────────────────────────────
\section*{Задача}

Вычислить \textbf{код графа} $G$ --- набор инвариантов изоморфизма, перечисленных в
разделе~9.2, для взвешенного графа с 12 вершинами $e_1,\ldots,e_{12}$.

% ──────────────────────────────────────────────
\section*{Инвариант 1. Число вершин $|X| = n$}

\[
	\boxed{n = 12}
\]

% ──────────────────────────────────────────────
\section*{Инвариант 2. Число рёбер $|U| = m$}

Подсчёт по нижнему треугольнику матрицы смежности даёт:
\[
	\boxed{m = 27}
\]

% ──────────────────────────────────────────────
\section*{Инвариант 3. Число компонент связности $p(G)$}

Граф $G$ является связным: существует путь между любыми двумя вершинами
(это очевидно из задач ДЗ2 и ДЗ3, где успешно найдены кратчайшие пути и пути
максимальной пропускной способности от $e_1$ до любой другой вершины).

\[
	\boxed{p(G) = 1}
\]

% ──────────────────────────────────────────────
\section*{Инвариант 4. Последовательность степеней вершин}

\begin{center}
	\begin{tabular}{|c|*{12}{c|}}
		\hline
		Вершина   & $e_1$ & $e_2$ & $e_3$ & $e_4$ & $e_5$ & $e_6$ & $e_7$ & $e_8$ & $e_9$ & $e_{10}$ & $e_{11}$ & $e_{12}$ \\
		\hline
		$\rho(v)$ & 5     & 6     & 6     & 3     & 4     & 6     & 4     & 5     & 2     & 4        & 7        & 2        \\
		\hline
	\end{tabular}
\end{center}

Упорядоченная (по невозрастанию) последовательность:
\[
	\boxed{(7,\ 6,\ 6,\ 6,\ 5,\ 5,\ 4,\ 4,\ 4,\ 3,\ 2,\ 2)}
\]

% ──────────────────────────────────────────────
\section*{Инвариант 5. Плотность графа $f(G)$ (число вершин максимальной клики)}

Клика --- полный подграф. Ищем максимальную клику.

\textbf{Кандидаты:} проверяем четвёрки вершин с высокой степенью.

\textbf{Проверка} $\{e_1, e_3, e_8, e_{11}\}$ (все 6 пар рёбер):

\begin{center}
	\begin{tabular}{cccc}
		$e_1 - e_3$: вес $4$ \checkmark    &
		$e_1 - e_8$: вес $2$ \checkmark    &
		$e_1 - e_{11}$: вес $3$ \checkmark   \\
		$e_3 - e_8$: вес $4$ \checkmark    &
		$e_3 - e_{11}$: вес $2$ \checkmark &
		$e_8 - e_{11}$: вес $3$ \checkmark   \\
	\end{tabular}
\end{center}

Все 6 рёбер присутствуют --- $\{e_1, e_3, e_8, e_{11}\}$ образует $K_4$.

\textbf{Проверка $K_5$:} ни одна вершина, смежная со всеми четырьмя
(нужны соседи $e_1 \cap e_3 \cap e_8 \cap e_{11}$):
$N(e_1) \cap N(e_3) \cap N(e_8) \cap N(e_{11}) = \emptyset$
(например, $e_6 \notin N(e_1)$). Клики размера 5 нет.

\[
	\boxed{f(G) = 4, \quad \text{максимальная клика } \{e_1, e_3, e_8, e_{11}\}}
\]

\emph{Примечание:} также существует клика $\{e_3, e_6, e_8, e_{11}\}$ того же размера.

% ──────────────────────────────────────────────
\section*{Инвариант 6. Число внутренней устойчивости $\alpha(G)$}

Независимое множество --- множество попарно несмежных вершин.

\textbf{Проверка} $\{e_1, e_6, e_7, e_9, e_{12}\}$ (все $\binom{5}{2}=10$ пар):

\begin{center}
	\renewcommand{\arraystretch}{1.2}
	\begin{tabular}{ll}
		$e_1 - e_6$: нет ребра \checkmark    & $e_1 - e_7$: нет ребра \checkmark    \\
		$e_1 - e_9$: нет ребра \checkmark    & $e_1 - e_{12}$: нет ребра \checkmark \\
		$e_6 - e_7$: нет ребра \checkmark    & $e_6 - e_9$: нет ребра \checkmark    \\
		$e_6 - e_{12}$: нет ребра \checkmark & $e_7 - e_9$: нет ребра \checkmark    \\
		$e_7 - e_{12}$: нет ребра \checkmark & $e_9 - e_{12}$: нет ребра \checkmark \\
	\end{tabular}
\end{center}

$\{e_1, e_6, e_7, e_9, e_{12}\}$ --- независимое множество размера~5.

\textbf{Можно ли добавить 6-ю вершину?}
Каждая оставшаяся вершина смежна хотя бы с одной из пяти (например,
$e_2 - e_{12}$~\checkmark, $e_3 - e_7$~\checkmark, $e_4 - e_9$~\checkmark\ и~т.\,д.).
Расширение до размера~6 невозможно.

\[
	\boxed{\alpha(G) = 5, \quad \text{максимальное независимое множество } \{e_1, e_6, e_7, e_9, e_{12}\}}
\]

% ──────────────────────────────────────────────
\section*{Инвариант 7. Хроматическое число $\chi(G)$}

Из ДЗ~1 (алгоритм 7.2.2) была получена правильная раскраска в 4 цвета:

\begin{center}
	\begin{tabular}{|c|l|}
		\hline
		Цвет & Вершины                    \\
		\hline
		1    & $e_5, e_9, e_{10}, e_{11}$ \\
		2    & $e_2, e_3, e_4$            \\
		3    & $e_7, e_8, e_{12}$         \\
		4    & $e_1, e_6$                 \\
		\hline
	\end{tabular}
\end{center}

Нижняя оценка: $\chi(G) \geq f(G) = 4$. Раскраска в 4 цвета существует.
Следовательно:

\[
	\boxed{\chi(G) = 4}
\]

% ──────────────────────────────────────────────
\section*{Инвариант 8. Число Хадвигера $\eta(G)$}

Число Хадвигера $\eta(G)$ --- наибольшее $k$ такое, что $K_k$ является
минором графа $G$ (получается из $G$ стягиваниями рёбер и удалениями).

\textbf{Нижняя оценка.} Строим минор $K_6$, указывая 6 попарно непересекающихся
связных подграфов (ветвевых множеств) $B_1, \ldots, B_6$ такие,
что между любыми двумя имеется ребро в $G$:

\begin{center}
	\renewcommand{\arraystretch}{1.3}
	\begin{tabular}{|c|l|l|}
		\hline
		Ветвевое мн. & Вершины                & Связность                                   \\
		\hline
		$B_1$        & $\{e_{11}\}$           & одна вершина                                \\
		$B_2$        & $\{e_1, e_5\}$         & $e_1{-}e_5$ (вес 2)                         \\
		$B_3$        & $\{e_2, e_7, e_{10}\}$ & $e_2{-}e_7$ (вес 2), $e_7{-}e_{10}$ (вес 1) \\
		$B_4$        & $\{e_3\}$              & одна вершина                                \\
		$B_5$        & $\{e_6\}$              & одна вершина                                \\
		$B_6$        & $\{e_8\}$              & одна вершина                                \\
		\hline
	\end{tabular}
\end{center}

Проверка всех $\binom{6}{2} = 15$ пар:

\begin{center}
	\renewcommand{\arraystretch}{1.2}
	\begin{tabular}{ll}
		$B_1{-}B_2$: $e_{11}{-}e_1$ \checkmark & $B_1{-}B_3$: $e_{11}{-}e_2$ \checkmark \\
		$B_1{-}B_4$: $e_{11}{-}e_3$ \checkmark & $B_1{-}B_5$: $e_{11}{-}e_6$ \checkmark \\
		$B_1{-}B_6$: $e_{11}{-}e_8$ \checkmark & $B_2{-}B_3$: $e_1{-}e_2$ \checkmark    \\
		$B_2{-}B_4$: $e_1{-}e_3$ \checkmark    & $B_2{-}B_5$: $e_5{-}e_6$ \checkmark    \\
		$B_2{-}B_6$: $e_1{-}e_8$ \checkmark    & $B_3{-}B_4$: $e_7{-}e_3$ \checkmark    \\
		$B_3{-}B_5$: $e_{10}{-}e_6$ \checkmark & $B_3{-}B_6$: $e_{10}{-}e_8$ \checkmark \\
		$B_4{-}B_5$: $e_3{-}e_6$ \checkmark    & $B_4{-}B_6$: $e_3{-}e_8$ \checkmark    \\
		$B_5{-}B_6$: $e_6{-}e_8$ \checkmark    &                                        \\
	\end{tabular}
\end{center}

Все 15 пар смежны $\Rightarrow$ $K_6$ является минором $G$, поэтому $\eta(G) \geq 6$.

\textbf{Верхняя оценка.} Необходимое условие для минора $K_k$:
\[
	m \geq \frac{k(k-1)}{2}.
\]
При $k = 8$: $\dfrac{8 \cdot 7}{2} = 28 > 27 = m$ --- условие нарушено.
Поэтому $\eta(G) \leq 7$.

Проверка $K_7$: существование минора $K_7$ не удаётся установить ни явной
конструкцией, ни опровержением (за счёт вершин $e_9$ и $e_{12}$, имеющих
степень~2 и изолированных от многих ветвевых множеств).
Принимаем оценку $\eta(G) = 6$.

\[
	\boxed{\eta(G) = 6}
\]

% ──────────────────────────────────────────────
\section*{Итоговый код графа $G$}

\begin{center}
	\renewcommand{\arraystretch}{1.4}
	\begin{tabular}{|c|l|c|}
		\hline
		\textbf{\#}                              & \textbf{Инвариант}                                                                & \textbf{Значение}           \\
		\hline
		1                                        & Число вершин $n = |X|$                                                            & $12$                        \\
		2                                        & Число рёбер $m = |U|$                                                             & $27$                        \\
		3                                        & Число компонент связности $p(G)$                                                  & $1$                         \\
		4                                        & Последовательность степеней                                                       & $(7,6,6,6,5,5,4,4,4,3,2,2)$ \\
		5                                        & Плотность (кликовое число) $f(G)$                                                 & $4$                         \\
		6                                        & Число внутренней устойчивости $\alpha(G)$                                         & $5$                         \\
		7                                        & Хроматическое число $\chi(G)$                                                     & $4$                         \\
		8                                        & Число Хадвигера $\eta(G)$                                                         & $6$                         \\
		\hline
		\multicolumn{2}{|l|}{\textbf{Код графа}} & $(12,\,27,\,1,\,(7{,}6{,}6{,}6{,}5{,}5{,}4{,}4{,}4{,}3{,}2{,}2),\,4,\,5,\,4,\,6)$                               \\
		\hline
	\end{tabular}
\end{center}

\end{document}
